\section{FORMAL RESTATEMENT}
\textbf{Erd\H{o}s \#27; independent reconstruction attempt, 2026-09-23.}
For distinct moduli $N\leq n_1<\cdots<n_t\leq CN$, choose one residue
class modulo each modulus. Let $\delta$ be the natural density of integers
in none of these classes. The question asks whether some fixed $C>1$
permits $\delta\leq\epsilon$ for every $\epsilon>0$ and every $N$.
All sets here are periodic modulo the least common multiple of the
moduli, so their natural densities exist.

\section{QUICK LITERATURE AND CONTEXT CHECK}
The current page marks the assertion disproved. The primary paper of
Filaseta, Ford, Konyagin, Pomerance and Yu, \emph{Sieving by large
integers and covering systems of congruences}, proves the needed
uniform estimate in Theorem B. Below I prove elementary bounds and the
paper's useful dependency-error lemma, then make the exact deduction
from Theorem B explicit. The full sieve argument behind that theorem
is an external input, not a proof supplied by this attempt.

\section{WORK}
\subsection{1. What direct counting proves}
The union bound gives
\[
 \delta\geq1-\sum_i\frac1{n_i}
 \geq1-\sum_{m=N}^{\lfloor CN\rfloor}\frac1m
 =1-\log C-o(1).
\]
Thus no $1<C<e$ works. This does not address arbitrary fixed $C$:
once the reciprocal sum exceeds one the bound can become negative.

Conversely, choose each residue independently and uniformly modulo its
modulus. For a fixed integer $x$, its probability of surviving all
classes is $\prod_i(1-1/n_i)$. Averaging over one common period and then
over choices of residues shows that some choice has
\[
 \delta\leq\alpha:=\prod_i(1-1/n_i).
\]
If every integer modulus from $N$ to $M$ is used, this product telescopes
to $(N-1)/M$. In a bounded-ratio interval the random benchmark is
therefore about $1/C$, rather than zero.

\subsection{2. An elementary lower bound including dependence}
Define
\[
 \beta=\sum_{\substack{i<j\\\gcd(n_i,n_j)>1}}\frac1{n_i n_j}.
\]
We prove $\delta\geq\alpha-\beta$, the elementary lemma used in the
primary paper. Suppose the first $t-1$ classes leave a set $R$ of
density $\delta'$. Discard from these classes those whose moduli are
not coprime to $n_t$, leaving a larger surviving set $R_0$. A union
bound gives
\[
 d(R_0)\leq\delta'+
       \sum_{\substack{i<t\\\gcd(n_i,n_t)>1}}\frac1{n_i}.
\]
The set $R_0$ has a period coprime to $n_t$. The Chinese remainder
theorem therefore shows that exactly the fraction $1/n_t$ of $R_0$
lies in the newly selected residue class. The part removed from $R$
is contained in that intersection. Hence
\[
 \delta\geq(1-1/n_t)\delta'
       -\sum_{\substack{i<t\\\gcd(n_i,n_t)>1}}\frac1{n_i n_t}.
\]
Induction, starting with one class, proves the claimed inequality:
multiplying the preceding error by $1-1/n_t$ can only reduce its size.
In particular, pairwise coprime moduli have the exact density $\alpha$,
independent of the selected residues.

\subsection{3. The verified external theorem and the original quantifiers}
Theorem B of the cited paper states, in particular, that for every fixed
$K>1$ and every set of distinct integer moduli $S\subset(X,KX]$, the
minimum possible uncovered density satisfies
\[
 \delta^-(S)=(1+o_{X\to\infty}(1))
       \prod_{m\in S}(1-1/m),                         \tag{1}
\]
uniformly over $S$. Its published range actually allows growing $K$;
only the fixed-$K$ case is needed here.

Fix any proposed $C>1$. For sufficiently large integer $N$, apply (1)
with $X=N-1$ and $K=2C$. Every allowed modulus is then in $(X,KX]$.
As every factor belongs to $(0,1)$, deleting moduli can only increase
the product, so
\[
 \prod_i(1-1/n_i)
 \geq\prod_{m=N}^{\lfloor CN\rfloor}(1-1/m)
 =\frac{N-1}{\lfloor CN\rfloor}
 \geq\frac1{2C}.
\]
By the uniformity in (1), for all sufficiently large $N$ every residue
system has $\delta\geq1/(4C)$. Choosing $\epsilon=1/(8C)$ contradicts
the required assertion. Thus the published theorem gives the negative
answer for every fixed $C$, with the endpoint $n_i=N$ handled explicitly.

\section{VERIFICATION AND REMAINING GAP}
The averaging estimate is an existence upper bound and cannot be
reversed merely by declaring the congruences independent. The proof of
the lower bound only invokes independence when the Chinese remainder
theorem justifies it. The error $\beta$ need not vanish for all moduli
in a fixed-ratio interval; consequently the elementary lemma alone
does not prove (1). Establishing the uniform sieve estimate in (1)
is the main part of the known proof still not reconstructed here.

\section{FINAL STATUS}
\textbf{UNRESOLVED (independent proof reconstruction incomplete).}
The problem itself is known to be disproved. This attempt fully proves
the range $1<C<e$, the averaging upper bound and the dependency-error
lemma, and verifies the deduction of the full answer from an explicitly
cited external theorem. It does not claim an independent complete proof
of that theorem or a new result.

\section{SOURCES}
Problem and status: \url{https://www.erdosproblems.com/27}.
M.~Filaseta, K.~Ford, S.~Konyagin, C.~Pomerance and G.~Yu,
\emph{Sieving by large integers and covering systems of congruences},
J.~Amer.~Math.~Soc.~20 (2007), 495--517, Lemma~2.1 and Theorem~B:
\url{https://www.math.kent.edu/~yu/research/covering.pdf}.
Both checked 2026-09-23.

This is a subjective estimate of the independent reconstruction,
not a probability that the already published result is true.

\section{COMPLETION ESTIMATE}
\textbf{COMPLETION: 30\%}
